\documentclass[11pt]{article}
\usepackage[margin=1in]{geometry}
\usepackage{amsmath,amssymb,amsthm,mathtools}
\usepackage{graphicx}
\usepackage[hidelinks]{hyperref}
\usepackage{microtype}

\newtheorem{theorem}{Theorem}
\newtheorem{lemma}[theorem]{Lemma}
\newtheorem{corollary}[theorem]{Corollary}
\newcommand{\R}{\mathbb R}
\newcommand{\D}{\mathrm D}
\newcommand{\Id}{\mathrm I}
\newcommand{\essran}{\overline{\operatorname{ran}}}

\title{Uniqueness for spatially homogeneous rank-one-cone diffusions}
\author{Candidate full sufficient-condition theorem for arXiv:2605.00439}
\date{11 August 2026}

\begin{document}
\maketitle

\begin{center}
\fbox{\parbox{0.92\linewidth}{\textbf{Status.}
Candidate full theorem, likely valid, pending independent expert review.
It gives weighted one-sided stability, comparison, and uniqueness for every
bounded datum whenever the coefficient is a nonnegative combination of
fixed rank-one directions with state-dependent weights.}}
\end{center}

\section{Source question and result}

Auscher and Bechtel consider bounded weak solutions on $\R^n$ of
\begin{equation}\label{eq:QL}
 \partial_tu-\operatorname{div}(a(t,x,u)\nabla u)=0,
 \qquad u(0)=u_0.
\end{equation}
Their Corollary~1.5 constructs a global bounded weak solution for every
bounded datum, but their concluding uniqueness problem asks for suitable
conditions ensuring uniqueness at this rough initial regularity
\cite{AB}.

\begin{center}
\includegraphics[width=0.98\linewidth]{figures/open_problem_crop.png}
\end{center}

Here is a broad structural answer.  Fix nonzero vectors
$e_1,\ldots,e_m\in\R^n$ spanning $\R^n$.  Assume
\begin{equation}\label{eq:cone}
 a(t,x,s)=a(t,s)=\sum_{r=1}^m\alpha_r(t,s)e_r\otimes e_r,
 \qquad \alpha_r(t,s)>0,
\end{equation}
where each $\alpha_r$ is measurable in $t$, locally Lipschitz in $s$, and,
on every bounded time interval and compact state interval, is bounded above
and away from zero.  Assume also that the resulting matrix $a$ satisfies
Assumption~A of the source.  Thus the weights may vary with time and the
state while the elliptic directions remain fixed.  Isotropic diffusion
$a(t,s)=\alpha(t,s)\Id$ and fixed-axis diagonal anisotropy are special cases.

For $p>n/2$, put
\begin{equation}\label{eq:weight}
 \omega_p(x)=(1+|x|^2)^{-p}.
\end{equation}

\begin{theorem}[Weighted comparison and uniqueness]\label{thm:main}
Assume \eqref{eq:cone}.  Let $u,v$ be bounded weak solutions of
\eqref{eq:QL}, in the precise sense of \cite[Definition 1.2]{AB}, with
bounded initial data $u_0,v_0$.  For every $T<\infty$ there is
$C_T<\infty$, depending only on $p$, the fixed directions, and the bounds
of the $\alpha_r$ on the joint solution range up to time $T$, such that
for almost every $t\in(0,T)$,
\begin{equation}\label{eq:stability}
 \int_{\R^n}(u(t)-v(t))_+\omega_p
 \leq e^{C_Tt}\int_{\R^n}(u_0-v_0)_+\omega_p.
\end{equation}
The analogous estimate holds for the negative part.  Consequently:
\begin{enumerate}
\item $u_0\leq v_0$ almost everywhere implies $u\leq v$ almost everywhere;
\item two bounded weak solutions with the same bounded datum coincide;
\item under \eqref{eq:cone}, the solution supplied by source
Corollary~1.5 is the unique global bounded weak solution.
\end{enumerate}
\end{theorem}

The argument has two ingredients.  The fixed rank-one structure converts
the difference equation into a sum of one-dimensional monotone diffusions.
A directional Kato inequality then propagates a polynomially weighted
$L^1$ distance.  The source's own local $L^2$ trace theorem closes the
estimate at $t=0$, even though the solutions need not be integrable on
$\R^n$.

\section{Directional Kato lemma}

Write $\D_e=e\cdot\nabla$.  We use the standard parabolic Kato inequality
in the following finite-direction form.

\begin{lemma}[Finite-direction Kato inequality]\label{lem:kato}
Let $w,z_1,\ldots,z_m$ be locally integrable on
$(0,T)\times\R^n$, with enough local energy regularity to justify the weak
equation
\begin{equation}\label{eq:abstract}
 \partial_tw=\sum_{r=1}^m\D_{e_r}^2z_r.
\end{equation}
Suppose $z_r=q_rw$ with $q_r\in H^1_{\mathrm{loc}}$ locally bounded and
strictly positive (locally bounded away from zero).
Then, in distributions,
\begin{equation}\label{eq:kato}
 \partial_t(w)_+\leq
 \sum_{r=1}^m\D_{e_r}^2(z_r)_+.
\end{equation}
The same assertion holds for negative parts and, after addition, for
absolute values.
\end{lemma}

\begin{proof}
For smooth $w,q_r$, the claim is transparent. Since $q_r>0$,
$(z_r)_+=q_rw_+$. Along a line parallel to $e_r$, the distributional chain
rule gives
\[
 \D_{e_r}^2(q_rw_+)
 -\mathbf1_{\{w>0\}}\D_{e_r}^2(q_rw)
 =q_r\,\delta_{\{w=0\}}|\D_{e_r}w|^2\geq0;
\]
all first and second derivatives of $q_r$ cancel. Also
$\partial_tw_+=\mathbf1_{\{w>0\}}\partial_tw$. Multiplying
\eqref{eq:abstract} by this common indicator and summing the displayed
inequalities proves \eqref{eq:kato}.

For the stated weak regularity, take Steklov averages in time and smooth
convex approximations $j_\varepsilon$ of $s\mapsto s_+$, perform the same
calculation after spatial mollification, and pass first in the mollifier,
then the Steklov scale, and finally $\varepsilon$. The
$H^1_{\mathrm{loc}}$ regularity of $q_r$, its local bounds, and the local
energy bounds control the commutators; the only surviving convexity term is
the nonnegative measure displayed above. This is the standard parabolic Kato
approximation, applied separately in each fixed direction and then summed.
\end{proof}

For the solutions considered below, $u,v$ and their spatial gradients are
locally square integrable, and every nonlinear primitive is locally
$H^1$ by the source's local Lipschitz assumption.  Thus Lemma~\ref{lem:kato}
applies by the displayed approximation argument.  This is also the familiar
$L^1$-accretivity mechanism for scalar nonlinear diffusion; compare the
bounded, non-integrable entropy theory in \cite{EJ}.

\section{Difference equation and weighted estimate}

Fix $s_*\in O$.  For each direction define the increasing primitive
\begin{equation}\label{eq:beta}
 \beta_r(t,s)=\int_{s_*}^s\alpha_r(t,q)\,dq.
\end{equation}
Since $\alpha_r$ is independent of $x$, the Sobolev chain rule gives
\[
 \operatorname{div}(a(t,u)\nabla u)
 =\sum_{r=1}^m\D_{e_r}^2\beta_r(t,u)
\]
in distributions.  Set
\[
 w=u-v,
 \qquad z_r=\beta_r(t,u)-\beta_r(t,v).
\]
More explicitly,
\[
 z_r=q_rw,\qquad
 q_r(t,x)=\int_0^1\alpha_r(t,v(t,x)+\theta w(t,x))\,d\theta.
\]
The local Lipschitz assumption and $u,v\in H^1_{\mathrm{loc}}$ give
$q_r\in H^1_{\mathrm{loc}}$ in space.
Subtracting the weak equations gives \eqref{eq:abstract}. Because every
$\alpha_r$ is positive,
\begin{equation}\label{eq:signbound}
 \operatorname{sgn}z_r=\operatorname{sgn}w\quad(z_r\ne0),
 \qquad |z_r|\leq L_{r,T}|w|,
\end{equation}
where $L_{r,T}$ is the supremum of $\alpha_r$ on the compact joint range
of $u,v$ and on $[0,T]$.  Lemma~\ref{lem:kato} therefore yields
\begin{equation}\label{eq:kato2}
 \partial_tw_+\leq\sum_r\D_{e_r}^2(z_r)_+.
\end{equation}

The weight \eqref{eq:weight} is integrable because $p>n/2$, and direct
differentiation gives
\begin{equation}\label{eq:hessian}
 |\D_e^2\omega_p(x)|
 \leq C_p|e|^2\omega_p(x),
 \qquad C_p=2p+4p(p+1).
\end{equation}
Test \eqref{eq:kato2} against increasingly large compactly supported
approximations to $\omega_p$.  The cutoff errors vanish by boundedness of
$z_r$ and integrability of $\omega_p$ and its first two derivatives.
Using \eqref{eq:signbound}--\eqref{eq:hessian}, in the sense of
distributions in time,
\begin{align}\label{eq:gronwall}
 \frac{d}{dt}\int w_+\omega_p
 &\leq\sum_r\int (z_r)_+\D_{e_r}^2\omega_p \\
 &\leq C_p\sum_rL_{r,T}|e_r|^2\int w_+\omega_p
 =C_T\int w_+\omega_p.
\end{align}

It remains only to identify the value at time zero.  A bounded quasilinear
solution is also a bounded weak solution of the linear equation with
coefficient $A(t,x)=a(t,u(t,x))$.  On the compact solution range, $A$ is
bounded and uniformly elliptic.  Hence \cite[Proposition 3.4(b)]{AB} gives
\begin{equation}\label{eq:l2trace}
 u(t)\to u_0,\qquad v(t)\to v_0
 \quad\hbox{in }L^2_{\mathrm{loc}}(\R^n)\quad(t\downarrow0).
\end{equation}
The weighted $L^1$ convergence follows: on a fixed ball use
\eqref{eq:l2trace} and Cauchy--Schwarz, while outside the ball use the
uniform bounds on $u,v,u_0,v_0$ and then send the radius to infinity because
$\omega_p\in L^1$.  Therefore
\[
 \int (u(t)-v(t))_+\omega_p
 \longrightarrow\int(u_0-v_0)_+\omega_p.
\]
Gronwall's lemma applied to \eqref{eq:gronwall} proves
\eqref{eq:stability}.  Reversing $u,v$ proves the negative-part estimate,
and all conclusions of Theorem~\ref{thm:main} follow.

\section{Examples, boundary of the method, and novelty audit}

The theorem includes
\[
 a(t,s)=\alpha(t,s)\Id,
 \qquad
 a(t,s)=\operatorname{diag}(\alpha_1(t,s),\ldots,\alpha_n(t,s)),
\]
and overcomplete fixed frames
$a(t,s)=\sum_r\alpha_r(t,s)e_r\otimes e_r$.  The weights need not coincide,
and may be time dependent.  The source permits local degeneracy as
$|s|\to\infty$; only uniform ellipticity on the bounded joint solution range
is needed.

The proof does not cover rotating eigendirections depending on $s$, nor
$x$-dependent coefficients: differentiating a state primitive then creates
spatial drift terms, and sign-separated directional Kato no longer applies
directly.  Those are genuine remaining parts of the source's broad
uniqueness program.

Exact-title, exact-equation, Kato, nonlinear-diffusion, and rank-one-cone
searches through 11 August 2026 found classical isotropic entropy and
$L^1$-contraction results, notably \cite{EJ}, but no statement matching the
fixed-rank-one cone for the source's bounded weak-solution class and rough
distributional trace.  The present argument is self-contained modulo the
standard Kato approximation detailed in Lemma~\ref{lem:kato} and the source's
own local $L^2$ trace theorem.

One final triage point is worth recording.  The source's VMO existence bullet
cannot literally enlarge its \emph{bounded} weak-solution class to unbounded
VMO traces: Proposition~3.4(b) already forces every bounded solution trace
to lie in $L^\infty$.  A separate supporting note in this packet's companion
attempt gives an explicit locally bounded unbounded CMO datum.  The
rank-one-cone uniqueness theorem above is the genuinely live result.

\begin{thebibliography}{9}
\bibitem{AB}
P.~Auscher and S.~Bechtel,
\emph{Existence and uniqueness of weak solutions to quasilinear PDEs with
critical data}, arXiv:2605.00439v3, 2026.

\bibitem{EJ}
J.~Endal and E.~R. Jakobsen,
\emph{$L^1$ contraction for bounded (non-integrable) solutions of degenerate
parabolic equations}, SIAM J. Math. Anal. \textbf{46} (2014), 3957--3982;
arXiv:1404.6418v2.
\end{thebibliography}

\end{document}
